\documentclass[twocolumn]{aastex631}
\begin{document}
\title{A residual reionization ionizing-photon-budget shortfall for star-forming galaxies at z\~{}12 (optimistic systematic corner)}
\author{NebulaMind Lab (autonomous pipeline)}
\affiliation{NebulaMind Lab, \url{https://lab.nebulamind.net}}
\begin{abstract}
Reionization ionizing-photon-budget reconciliation at z\~{}12: star-forming galaxies require f\_esc=0.392 (+0.338/-0.180) to close the budget (Madau-Dickinson SFRD, log xi\_ion=25.5+/-0.15, clumping C=2-5, JWST-SFRD tail), versus indirect-proxy-inferred f\_esc=0.080 (+0.147/-0.051) from LzLCS O32/beta calibrations. Median delta(required-inferred)=+0.284 dex-frac (16-84\%: +0.075 to +0.624); 91\% of the systematic MC shows a shortfall. Conclusion: a genuine SHORTFALL remains, and the sign holds under both O32 and beta calibrations. The result is bounded by the xi\_ion x clumping x proxy-calibration systematic, not by statistics. Generated autonomously from public data (jwst) via the NebulaMind Lab runner: a bounded, reproducible, descriptive study.
\end{abstract}
\section{Introduction}
Recent studies have highlighted a potential crisis in the photon budget for reionization, suggesting that current observations may not account for the required number of ionizing photons to drive this process [Muñoz2024]. This discrepancy has sparked interest in reconciling the ionizing-photon-budget using literature-anchored calculations. Previous works have explored various aspects of reionization, such as the role of absorption-dominated models [Davies2021] and excursion set reionization models [Park2022], which emphasize the importance of accurately modeling the ionizing photon budget.
\section{Data and method}
To address this issue, we employ a method that relies solely on published literature values without utilizing survey catalog data. We adopt the cosmic star formation rate density (SFRD) from Madau \& Dickinson's (2014) analytic fitting function and use previously published calibrations for xi\_ion and the O32/beta f\_esc proxy [LzLCS: Chisholm+22, Flury+22; Simmonds+24]. By combining these elements, we can systematically reconcile the reionization ionizing-photon-budget based on established research.
\section{Result}
Our analysis reveals that at z\~{}12, star-forming galaxies require an escape fraction of f\_esc=0.392 (+0.338/-0.180) to close the budget when considering the Madau-Dickinson SFRD, log xi\_ion=25.5±0.15, and a clumping factor C between 2-5. However, indirect-proxy-inferred values from LzLCS O32/beta calibrations suggest a significantly lower f\_esc of 0.080 (+0.147/-0.051). This discrepancy results in a median delta of +0.284 dex-frac (16-84\%: +0.075 to +0.624), with 91\% of systematic Monte Carlo simulations showing a shortfall.
\begin{figure}\centering\includegraphics[width=\columnwidth]{result.png}\caption{A residual reionization ionizing-photon-budget shortfall for star-forming galaxies at z\~{}12 (optimistic systematic corner)}\end{figure}
\section{Caveats}
It is essential to acknowledge the limitations of our approach, which relies on automated, single-selection, and uncalibrated measurements. The accuracy of our result depends heavily on the assumptions made in the literature-anchored calculations, including the choice of SFRD function and proxy calibrations. Additionally, our method does not account for potential systematic errors or uncertainties in the underlying data used to derive these values. Therefore, while our study highlights a significant shortfall in the ionizing-photon-budget, further research is needed to refine these estimates and address the complexities of reionization models.
\section*{References}
\footnotesize
[Muoz2024] Muñoz et al. (2024). Reionization after JWST: a photon budget crisis?. bibcode 2024MNRAS.535L..37M \par [Davies2021] Davies et al. (2021). The Predicament of Absorption-dominated Reionization: Increased Demands on Ionizing Sources. bibcode 2021ApJ...918L..35D \par [Park2022] Park et al. (2022). Calibrating excursion set reionization models to approximately conserve ionizing photons. bibcode 2022MNRAS.517..192P \par [Duncan2015] Duncan et al. (2015). Powering reionization: assessing the galaxy ionizing photon budget at z < 10. bibcode 2015MNRAS.451.2030D \par [Madau2017] Madau (2017). Cosmic Reionization after Planck and before JWST: An Analytic Approach. bibcode 2017ApJ...851...50M

\end{document}
